% ===========================================================================
%  sviluppatore/06-dashboard/10-view.tex
% ===========================================================================
\section{La View a tutto grafico}\label{sez:view-guscio}

La pagina View è il grafico: \file{templates/visual.html} ha un solo palco,
\cod{\#vd-main}, che il grafico riempie da bordo a bordo; tutto il resto gli
sta sopra in posizione assoluta. Il canvas dell'app perde margini e titolo
solo su questa pagina (\cod{\#qs-canvas[data-tab="visual"]} in
\file{static/css/visual.css}): l'altezza è quella della finestra meno
\cod{--ui-shell-inset}, il titolo resta nell'intestazione ritagliata a un
pixel per la sintesi vocale. Sotto i 640 pixel l'intestazione torna, perché lì
la navigazione è un menu e il pulsante che lo apre sta in lei.

\subsection{Il guscio}

\file{static/js/visual\_chrome.js}, con \simbolo{LyraChrome}\cod{.create(o)},
governa i tre pannelli che stanno sopra il grafico:

\begin{center}
\begin{tabular}{@{}L{34mm}L{104mm}@{}}
\toprule
\textbf{Pannello} & \textbf{Come si riduce} \\
\midrule
\rigaTenue
\cod{\#vd-topbar} & \cod{\#vd-top-fold}, il suo primo pulsante, la riduce alla
  linguetta \cod{\#vd-top-tab} nello stesso punto. \\
\cod{\#vd-toolbar} & \cod{\#vd-tools-fold} e \cod{\#vd-tools-tab}, allo stesso
  modo. \\
\rigaTenue
\cod{\#vd-ind-editor} & Un pannello agganciato in basso: \cod{\#vd-ind-min} lo
  riduce alla pillola \cod{\#vd-ind-pill}, in basso a sinistra, che lo
  riapre; \cod{\#vd-ind-max} lo massimizza e lo ripristina. \\
\bottomrule
\end{tabular}
\end{center}

Lo stato delle barre sta in due attributi del palco, \cod{data-top} e
\cod{data-tools}, che valgono \cod{"open"} o \cod{"folded"}: il CSS ne ricava
che cosa si vede. Le barre ridotte sfumano con \cod{visibility}, così escono anche
dall'ordine del tasto Tab. Ridurre sposta il fuoco sulla linguetta, riaprire
sul pulsante; la linguetta sta dove stava il pulsante, quindi il mouse resta
sopra al comando che serve dopo. Ridurre la barra in alto chiude i menu
aperti, con \cod{o.onFold}.

L'editor degli indicatori è un pannello in basso, a tutta larghezza, come su
TradingView. Il suo stato è l'attributo \cod{data-editor} del palco:
\cod{"closed"}, \cod{"open"} o \cod{"max"}. Aperto, il grafico finisce dove
il pannello comincia (\cod{bottom: var(--vd-ed-off)}): asse dei tempi,
riquadri delle anteprime e logo restano scoperti, e ciò che stava in basso
(intervallo, barra del replay, operazioni, stato) sale con lui. Massimizzato,
il pannello arriva sotto la barra in alto (\cod{--vd-ed-top}) e il grafico
resta della misura che aveva, sotto di lui.

\begin{perche}{Il grafico si accorcia, non si copre}
Un pannello sopra il grafico ne nasconde proprio la parte bassa, dove
compare il riquadro dell'indicatore che si sta scrivendo e dove sta l'asse
dei tempi. Accorciare il grafico costa una ridisegnata di Lightweight Charts
(\cod{autoSize}) e lascia leggibile l'anteprima.
\end{perche}

L'altezza (\cod{--vd-ed-h}) si cambia dal bordo superiore del pannello,
\cod{\#vd-ind-grip}, che per la tastiera e gli screen reader è un separatore
(\cod{role="separator"}, con \cod{aria-valuenow}): si trascina coi pointer
event e \cod{setPointerCapture}, o si sposta con le frecce (24 pixel, 96 con
Maiuscolo), Home ed End. \simbolo{limitiEditor} la tiene fra 170
pixel e quanto lascia al grafico 200 pixel; la prima volta vale il 42\%
dell'altezza del palco. Dallo stato massimizzato il trascinamento riparte
dall'altezza vera. Il doppio clic sull'intestazione (non sui suoi comandi) o
sul bordo massimizza e ripristina. All'apertura \cod{o.onEditor(true)}
carica il modello se l'editor è vuoto, e il fuoco va nel codice; riducendo,
va alla pillola.

Tutto resta in \cod{localStorage["vd-chrome"]}: le barre e l'editor,
\cod{\{open, h, max\}}. Una posizione di finestra salvata dalla versione di
prima (ha \cod{x}) si ignora. Uno script in linea subito dopo l'apertura del
palco rimette le barre ridotte prima del primo disegno, così non si vedono
aprirsi e richiudersi.

\subsection{L'impaginazione}

Dove stanno le cose che il grafico si porta sopra lo decide
\simbolo{impagina}, che scrive variabili CSS sul palco, e solo quando
cambiano:

\begin{center}
\begin{tabular}{@{}L{40mm}L{98mm}@{}}
\toprule
\textbf{Variabile} & \textbf{Che cosa} \\
\midrule
\rigaTenue
\cod{--vd-scale-w}, \cod{--vd-time-h} & La scala dei prezzi e l'asse dei
  tempi: nessun pannello ci va sopra. \\
\cod{--vd-rail-top}, \cod{--vd-rail-maxh} & Gli strumenti: sotto la barra in
  alto (anche quando va a capo), fino al logo di TradingView. \\
\rigaTenue
\cod{--vd-legend-top}, \cod{--vd-legend-left}, \cod{--vd-legend-maxw} & La legenda: nell'angolo fra le
  due barre o le loro linguette, larga fino alle etichette Sell/Buy o alla
  scala. \\
\cod{--vd-tags-top} & Sell/Buy: in alto, o sotto la barra se lei arriva fin
  lì. \\
\rigaTenue
\cod{--vd-panel-top} & Il pannello dell'ingranaggio, sotto Sell/Buy. \\
\cod{--vd-trades-top} & Le operazioni del replay, sotto la legenda. \\
\rigaTenue
\cod{--vd-dock-left} & La pillola e lo stato, dopo il logo. \\
\cod{--vd-ed-h}, \cod{--vd-ed-off}, \cod{--vd-ed-top} & L'editor: la sua
  altezza, dove finisce il grafico, dove arriva massimizzato. \\
\bottomrule
\end{tabular}
\end{center}

Le misure del grafico gliele dà \file{visual.js} con \simbolo{chartMetrics}:
la larghezza della scala dei prezzi, l'altezza dell'asse dei tempi e il
rettangolo del logo di TradingView, che non va mai coperto perché è
l'attribuzione chiesta dalla licenza. I pannelli si misurano con
\cod{offsetLeft} e \cod{offsetWidth}, non con \cod{getBoundingClientRect}:
durante la transizione sono scalati, e il rettangolo trasformato darebbe
posizioni sbagliate fino all'impaginazione successiva. Si impagina con un
\cod{ResizeObserver} (palco, barre, legenda, etichette, pillola), al
ridimensionamento della finestra e quando cambia la misura dell'asse dei
tempi, con \cod{subscribeSizeChange}, che scatta anche quando si allarga la
scala dei prezzi. \cod{o.onLayout(g)} passa a \file{visual.js} i rettangoli
del guscio sul lato sinistro, \cod{g.ostacoli}, solo quando cambiano;
\simbolo{layoutPaneLegends} sposta la legenda di un riquadro a destra di
quelli che la coprirebbero. Quando la colonna degli strumenti non ci sta
fra la barra in alto e il logo (con l'editor aperto, o in una finestra
bassa), \simbolo{impagina} la mette su due colonne (\cod{.vd-rail-2}),
misurando l'altezza di una colonna sola prima di decidere.

\paragraph{Sovrapposizione.} \cod{\#vd-chart} ha \cod{z-index: 0}. Senza, il
contenitore del grafico non è un contesto di sovrapposizione, e gli elementi
interni della libreria con uno z-index (i separatori dei riquadri, che si
trascinano) gareggiano con quelli del palco: passavano sopra i pannelli, e un
clic sul bordo di un riquadro sotto uno di loro arrivava al grafico. Il palco ha \cod{data-ui-group="none"}: \file{ui-system.js} prenderebbe
le due linguette per una serie di pulsanti e darebbe loro il raggio piccolo
dei gruppi.

\subsection{Le prove}

La sezione H di \file{tests/e2e\_grafici.js}: il grafico che riempie l'area e
la pagina che non scorre, il titolo, le barre e la legenda al loro posto, le
legende dei riquadri che scansano gli strumenti, il logo scoperto; ridurre e
riaprire (la linguetta sotto il mouse, il fuoco, la legenda che si sposta, il
grafico che non cambia misura, lo stato dopo il ricaricamento); il pannello
dell'editor (apertura dalla pillola, larghezza piena e altezza iniziale,
fuoco nel codice, grafico che finisce dove lui comincia, logo e asse scoperti,
strumenti sopra il logo, trascinamento e limiti, tastiera, massimizza e
ripristina, doppio clic, riduzione e riapertura alla stessa altezza,
ricaricamento); lo schermo stretto.
